\problemname{Long Staircases}

It is said that high up in the Alps live $C$ sisters, each of whom has her own favorite number $f$.
Every sister has therefore built a house $f$ units of height above the surface of the water.
The houses are built on a grid of size $R \times C$, such that the bottom left cell is $(0,0)$ and the top right cell is $(C-1,R-1)$.
The first sister has built her house at $(0,0)$, the second at $(1,0)$, and so on, so that the last sister has built her house at $(C-1,0)$.
In other words, the houses are built along the bottom row of the grid.

With the houses located so high up, it becomes difficult to get up and down.
Every sister has therefore built a staircase from her house up through the $R-1$ rows above it, where the height decreases by $1$ for every step.
So for every cell that is not on the bottom row, there is a cell below it whose height is $1$ greater.
For example, if a sister's house is at $(2,0)$ and has height $10$, she builds a step at $(2,1)$ with height $9$, a step at $(2,2)$ with height $8$, and so on.

You are currently on this grid.
More precisely, you are at the coordinates $(s_x, s_y)$ and want to get to $(g_x, g_y)$.
You can move to cells that are adjacent either up, left, down or right.
Going right or left is never a problem, even if there may be large differences in height.
Naturally, you cannot move to a coordinate outside the grid.

High altitudes are quite strenuous, so your goal is to minimize the time you spend at high altitude.
More precisely, if your walk from $(s_x, s_y)$ to $(g_x, g_y)$ visits a certain set of cells, we define the strain of the walk as the sum of the heights of all the cells you visit.
If you walk optimally, what is the smallest possible strain?

\begin{figure}[h]
  \centering
  \includegraphics[width=0.3\textwidth]{sample1.pdf}
    \\The picture shows what the grid looks like in sample $1$. The bottom row contains the sisters' houses.
    The start and the goal are also marked, along with the least straining walk between them.
\end{figure}

\section*{Input}
The first line of input contains $R$ and $C$ ($1 \leq R \leq 10^9$, $1 \leq C \leq 10$),
the dimensions of the grid.

The second line contains $C$ integers $f_1, \dots, f_C$ ($R-1 \leq f_i \leq 10^9$),
the height at which each sister's house is located.
$f_1$ is the height of the house at $(0,0)$, $f_2$ of the house at $(1,0)$, and so on.
\textbf{Note} that the heights $f_i$ are chosen so that the remaining heights in the grid are not negative.

The third line contains the integers $s_x$ and $s_y$ ($0 \leq s_x < C$, $0 \leq s_y < R$),
the coordinate where you start.

The fourth line contains the integers $g_x$ and $g_y$ ($0 \leq g_x < C$, $0 \leq g_y < R$),
the coordinate you want to get to.

\section*{Output}
Print an integer: the strain of the least straining walk from the start to the coordinate you want to get to.
\textbf{Note: the answer does not necessarily fit in a 32-bit integer.}

\section*{Scoring}
Your solution will be tested on a set of test groups.
To earn points for a group, you must pass all test cases in that group.

\noindent
\begin{tabular}{| l | l | p{12cm} |}
  \hline
  \textbf{Group} & \textbf{Points} & \textbf{Constraints} \\ \hline
  $1$    & $20$  & $C = 2$ \\ \hline
  $2$    & $40$  & $R \leq 10\,000$ \\ \hline
  $3$    & $40$  & No additional constraints. \\ \hline
\end{tabular}

\section*{Explanation of Sample 1}
Sample $1$ is the grid shown in the picture above.
The least straining walk goes right, right, up, up, up, up, left, left.
The strain is $15+12+5+4+3+2+1+8+11=61$.
